\documentclass[10pt,a4paper]{amsart}
\usepackage[margin=19mm]{geometry}
\usepackage{amsmath,amssymb,mathtools}
\usepackage[scr=rsfs,cal=euler]{mathalfa}
\usepackage{xfrac}
\usepackage[unicode,hidelinks,bookmarks=false]{hyperref}
\newcommand{\ie}{i.e.\ }
\newcommand{\cf}{cf.\ }
\newcommand{\resp}{respectively\ }
\newcommand{\Subsection}{\subsection}
\providecommand{\nobreakdash}{\nobreak\hbox{-}}
\setlength{\emergencystretch}{2em}
\multlinegap0pt
\usepackage{bm}
\usepackage{bbm}
\usepackage{dsfont}
\usepackage{xspace}
\usepackage{cancel}
\usepackage{mathtools}
\usepackage{amsthm}
\makeatletter
\@ifundefined{algorithm}{\newtheorem{algorithm}{Algorithm}}{}
\@ifundefined{proposition}{\newtheorem{proposition}[algorithm]{Proposition}}{}
\makeatother
\newcommand{\ZZ}{\mathbb{Z}}
\newcommand{\cO}{\mathcal{O}} 
\title[On the Squarefree Values of Degree-$2q$ Polynomials]{On the Squarefree Values of Degree-$2q$ Polynomials}
\author{Sergio Ricardo Zapata Ceballos, Fatemeh Jalalvand}
\date{Source: arXiv:2608.10335v1 (CC BY 4.0)}
\begin{document}
\maketitle

\noindent\emph{Source and licence.} On the Squarefree Values of Degree-$2q$ Polynomials. Version arXiv:2608.10335v1 (CC BY 4.0), distributed under the Creative Commons Attribution 4.0 International licence, as recorded in the arXiv licence metadata for this exact version and shown on the abstract page https://arxiv.org/abs/2608.10335v1. Full rights notice: licenses/arxiv-2608.10335-CC-BY-4.0.txt.\par\smallskip

\noindent\textit{Editorial scope.} Counted unit: the Proposition whose source label is \texttt{prop: sieve} in the article (source0.tex lines 280--292, source label \texttt{prop: sieve}), delivered together with the article's own complete original proof of it (source0.tex lines 294--389, one \texttt{proof} environment of 96 lines). No mathematics is written, repaired or re-derived: statement and proof are the article's own text line for line. Required context retained uncounted: prop: conjugatenotdividing (lines 258--261). Every definition, display and label of those lines is retained unchanged. Omitted from this package and not redistributed: 1--257, 262--279, 293--293, 390--639. Nothing inside a retained excerpt is omitted or altered: every retained body is delivered exactly as the article's lines are written, so no omission receipt of any kind accompanies this package. Citations: the retained excerpts carry no citation command at all, so no bibliography entry of the article is transcribed and no reference is redistributed. Reference adaptation: none was needed and none was made; every reference inside the delivered unit resolves inside it, so no unresolved reference or citation is produced. Printed slips kept literal: no printed word, symbol, hypothesis or number of the counted statement or of its proof was altered. Layout and class: the article's own class is \texttt{amsart} with packages including verbatim, xcolor, upgreek, amssymb, amsmath, mathtools, amsfonts, amsthm, eucal, fullpage, color, caption, subcaption, enumitem; the wrapper is the lane's pinned \texttt{amsart} editorial wrapper at 10pt on A4, which restates the theorem-like environments the retained text needs and adds the article's own macro definitions that the retained text needs (\texttt{\textbackslash ZZ}, \texttt{\textbackslash cO}), with extra wrapper packages (bm, bbm, dsfont, xspace, cancel, mathtools). The delivered numbering is the unit's own. Assets: no retained span contains a figure environment, a graphics-inclusion command, a PDF-page inclusion, a table or a TikZ picture; no image file is copied into this package, the package declares no asset dependency and no image file is redistributed with it. Licence: version arXiv:2608.10335v1 of this article is distributed under the Creative Commons Attribution 4.0 International licence, as recorded in the arXiv licence metadata for this exact version and shown on the abstract page https://arxiv.org/abs/2608.10335v1. A full rights notice is delivered at licenses/arxiv-2608.10335-CC-BY-4.0.txt. Characters: every retained excerpt is pure ASCII, so the delivered engine, XeLaTeX with its default Latin Modern text font, reports no missing character.
\begin{proposition}
\label{prop: conjugatenotdividing}
      Let \(\mathfrak{p}\) be a prime in $\cO_K$ and $n \in \ZZ$, with $p=\mathfrak{p}\cap \mathbb{Z}$  such that $\mathfrak{p} \mid (N_{L/K}(n - \alpha))$ and $p\nmid Disc(h)$. If $\bar{\mathfrak{p}}$ is a conjugate of $\mathfrak p$ and $\bar{\mathfrak{p}} \neq \mathfrak{p} $, then  $\bar{\mathfrak{p}} \nmid (N_{L/K}(n - \alpha))$.
\end{proposition}

\begin{proposition}
\label{prop: sieve}
Suppose that $\gcd\{h(n):n\in\mathbb{Z}\}$ is squarefree. Then, for each $p\in\mathcal{I}$, there exists an integer $r_p$ such that $h(r_p)\not\equiv 0\pmod{p^2}$. Moreover, there is a positive proportion of integers $n$ satisfying both
    \[
        n\equiv r_p\pmod{p^2\mathcal{O}_K}
        \qquad\text{for all }p\in\mathcal{I},
    \]
and
    \[
        \bigl(N_{L/K}(n-\alpha)\bigr)
    \]
is squarefree away from $\mathcal{I}_K$.
\end{proposition}

\begin{proof}
    For each prime ideal $\mathfrak{p}$ of $\mathcal{O}_K$, define
    \[
        \rho(\mathfrak{p}^2) \coloneqq \#\left\{\overline{n}\in\mathcal{O}_K/\mathfrak{p}^2
        \vcentcolon n\in\mathbb{Z} \text{ and } m(n)\equiv 0\pmod{\mathfrak{p}^2}\right\}.
    \]
    We first estimate $\rho(\mathfrak{p}^2)$. Suppose that
    $\mathfrak{p}\nmid 2\operatorname{Disc}(m)$. Then $m(x)$ has either no roots or two distinct roots modulo $\mathfrak{p}$. By Hensel's lemma, each root modulo $\mathfrak{p}$ lifts uniquely to a root modulo $\mathfrak{p}^2$. Hence
    \[
        \rho(\mathfrak{p}^2)\leq 2 
    \]
    for every $\mathfrak{p}\nmid 2\operatorname{Disc}(m)$. On the other hand, if
    $\mathfrak{p}\mid 2\operatorname{Disc}(m)$, then
    \[
        \rho(\mathfrak{p}^2)<p^2,
    \]
    where $p$ is the rational prime below $\mathfrak{p}$. Indeed, otherwise $m(n)\equiv0\pmod{\mathfrak{p}^2}$ for every $n\in\mathbb{Z}$, and hence
    \[
        h(n) = N_{K/\mathbb{Q}}(m(n)) \equiv0\pmod{p^2}
    \]
    for every $n\in\mathbb{Z}$, contradicting the assumption that $\gcd\{h(n):n\in\mathbb{Z}\}$ is squarefree.
    
    Fix $X\geq1$, and define
    \[
        \mathcal{P}_X \coloneqq \left\{\mathfrak{p}\notin\mathcal{I}_K \vcentcolon N_{K/\mathbb{Q}}(\mathfrak{p}^2)\leq X \right\}.
    \]
    Set
    \[ 
        P_X\coloneqq \prod_{\mathfrak{p}\in\mathcal{P}_X}\mathfrak{p}^2
        \qquad\text{and}\qquad M\coloneqq \prod_{p\in\mathcal{I}}p^2\mathcal{O}_K.
    \]
    Finally, let
    \[
        Z_X \coloneqq \left\{\overline{n}\in\mathcal{O}_K/MP_X \vcentcolon n\in\mathbb{Z}\right\}.
    \]
    
    Since $\gcd\{h(n):n\in\mathbb{Z}\}$ is squarefree, for each $p\in\mathcal{I}$ there exists an integer $r_p$ such that
    \[
        h(r_p)\not\equiv0\pmod{p^2}.
    \]
    
    For each $\mathfrak{p}\in\mathcal{P}_X$, define
    \[
        N_{\mathfrak{p}} \coloneqq \left\{\overline{n}\in\mathcal{O}_K/\mathfrak{p}^2 \vcentcolon n\in\mathbb{Z} \text{ and } m(n)\not\equiv0\pmod{\mathfrak{p}^2}\right\}.
    \]
    We then define
    \[
        N_X \coloneqq \left\{\overline{n}\in Z_X \vcentcolon
        \begin{array}{l}
            n\equiv r_p\pmod{p^2\mathcal{O}_K}
            \text{ for all }p\in\mathcal{I},\\[2mm]
            m(n)\not\equiv0\pmod{\mathfrak{p}^2}
            \text{ for all }\mathfrak{p}\in\mathcal{P}_X
        \end{array}
        \right\}.
    \]
    
    By the Chinese remainder theorem, the natural ring homomorphism
    \[
        \phi: \mathcal{O}_K/MP_X \longrightarrow \prod_{p\in\mathcal{I}} \mathcal{O}_K/p^2\mathcal{O}_K \times \prod_{\mathfrak{p}\in\mathcal{P}_X} \mathcal{O}_K/\mathfrak{p}^2
    \]
    is an isomorphism. By Proposition~\ref{prop: conjugatenotdividing}, the prime ideals occurring in the product are pairwise non-conjugate, and hence
    \[
        \phi(N_X) = \prod_{p\in\mathcal{I}}\{\overline{r_p}\} \times \prod_{\mathfrak{p}\in\mathcal{P}_X} N_{\mathfrak{p}}.
    \]
    Therefore,
    \[
        \#N_X = \prod_{\mathfrak{p}\in\mathcal{P}_X} \left(p^2-\rho(\mathfrak{p}^2)\right) = \left(\prod_{\mathfrak{p}\in\mathcal{P}_X}p^2\right) \prod_{\mathfrak{p}\in\mathcal{P}_X} \left(1-\frac{\rho(\mathfrak{p}^2)}{p^2}\right),
    \]
    where $p$ denotes the rational prime below $\mathfrak{p}$. Since every $\mathfrak{p}\in\mathcal{P}_X$ is unramified, we have
    \[ 
        \#\left\{\overline{n}\in\mathcal{O}_K/\mathfrak{p}^2 \vcentcolon n\in\mathbb{Z}\right\} = p^2.
    \]
    Consequently,
    \[  
        \#Z_X =\left(\prod_{p\in\mathcal{I}}p^2\right)\left(\prod_{\mathfrak{p}\in\mathcal{P}_X}p^2\right).
    \]
    It follows that
    \[
        \frac{\#N_X}{\#Z_X} = \frac{1}{\displaystyle\prod_{p\in\mathcal{I}}p^2} \prod_{\mathfrak{p}\in\mathcal{P}_X} \left(1-\frac{\rho(\mathfrak{p}^2)}{p^2}\right).
    \]
    
    For finitely many primes $\mathfrak{p}$, we have $\rho(\mathfrak{p}^2)\leq p^2$. For the other primes $\mathfrak{p}$, we have $\rho(\mathfrak{p}^2)\leq2$. Moreover, the corresponding rational primes $p$ are distinct by Proposition~\ref{prop: conjugatenotdividing}. Hence
    \[
        \sum_{\mathfrak{p} \in \mathcal{P}_X} \frac{\rho(\mathfrak{p}^2)}{p^2}
    \]
    converges, and therefore the infinite product
    \[
        \prod_{\mathfrak{p} \in \mathcal{P}_X}\left(1-\frac{\rho(\mathfrak{p}^2)}{p^2}\right)
    \]
    converges to a positive real number. Thus,
    \[
        \lim_{X\to\infty}\frac{\#N_X}{\#Z_X}>0.
    \]
    This proves the proposition.
\end{proof}

\end{document}
